\documentclass[a4paper]{article}

\usepackage{graphicx}
\usepackage{amsmath,amssymb}
\usepackage{minted}
\usepackage{multirow}
\usepackage{float}
\usepackage{todonotes}
\usepackage{forest}
\usepackage{xstring}
\usepackage{xcolor}
\usepackage[most]{tcolorbox}
\usepackage{xparse}
\usepackage{natbib}

\usetikzlibrary{graphs,graphdrawing}
\usegdlibrary{layered}

% for external graphviz
\input{/home/donkarlo/Dropbox/repo/nd_language_ai_project/src/nd_language_ai/formal/computer/programming/tex/latex/code_clips/external_graphviz.tex}

% for tags
\input{/home/donkarlo/Dropbox/repo/nd_language_ai_project/src/nd_language_ai/formal/computer/programming/tex/latex/parts/control_sequence/kind/semtag/semtag.tex}

% for links
\input{/home/donkarlo/Dropbox/repo/nd_language_ai_project/src/nd_language_ai/formal/computer/programming/tex/latex/code_clips/seehere.tex}

\title{Language contribuition in Self-awareness}
\date{}
\author{Mohammad Rahmani}

\begin{document}
    \maketitle
    
    
    \section{I am aware of my awareness}
        Language, particularly inner speech, is hypothesised to generate and sustain self-awareness or higher-order consciousness. The paper proposes the HOLISTIC model to integrate the neurobiology of language, inner speech, and consciousness. A core network for inner-speech production generates internal words and interacts with distributed sensory, motor, and emotional representations. Sustained conscious processing is proposed to involve the default mode network together with prefrontal, thalamic, and brainstem systems. The paper reviews converging evidence for this hypothesis, but does not experimentally prove that language is necessary or sufficient for self-awareness \cite{skipper-2022-a-voice-without-a-mouth-no-more-the-neurobiology-of-language-and-consciousness}.
        
        \cite{edelman-2011-biology-of-consciousness} integrates the Dynamic Core and Global Workspace theories into a biological account of consciousness. It distinguishes primary consciousness from higher-order consciousness, with primary consciousness not requiring linguistic capability. The authors argue that true language with syntax enables humans to develop higher-order consciousness. Language allows internal linguistic tokens to represent experiences beyond the immediate remembered present and enables awareness of one's own awareness. Higher-order consciousness consequently supports narrative thought about the past and future, while inner speech further elaborates this self-reflective capacity.
        
        
        
        \bibliographystyle{plainnat}
        \bibliography{/home/donkarlo/Dropbox/repo/nd_shared_research/bibliography/bibliography.bib}
\end{document}